\newchapter{Custom Commands}

\comment{ADD THIS SECTION.}

\newsection{Text Formatting}

\comment{ADD THIS SECTION.}

\newsubsection{Custom Date Formats}

\comment{ADD THIS SECTION.}

\newsubsection{Adding Custom Text Formats}

\comment{ADD THIS SECTION.}

\newsection{General Commands}

\comment{ADD THIS SECTION.}

\newsubsection{Figures and Tables}

\comment{ADD THIS SECTION.}

\newsubsection{Academic Environments}

\comment{ADD THIS SECTION.}

\newsubsection{Environment References}

\comment{ADD THIS SECTION.}

\newsubsection{Citations}

For this template, I have chosen to setup my citation commands inline with Harvard-style citation guidelines. 
As such, I have provided are two different formats for inserting a citation, direct and indirect. 
Their use-cases and implementations are described below.

\begin{command}{Direct Citation}
    A \underline{direct citation} is used when referring to a bibliography entry directly, for example when saying \say{according to \textit{bibliography\_entry}...}.
    Given the Harvard-style guidelines, this follows the general format:
    \begin{quote}
        \centering\commandformat{author, author, ..., and author (year\_published)}
    \end{quote}

    \noindent To call this function, we use the syntax given by \Cref{snippet:Direct citation format.}.

    \begin{newsnippet}{Direct citation format.}
        \begin{minted}{latex}
\directcitation{bibliography_entry_key}
        \end{minted}
    \end{newsnippet}

    \noindent The parameters for this function are shown below in \Cref{table:Direct citation parameters.}.

    \begin{newtable}{Direct citation parameters.}{|c|c|c|}
        \hline
        \textbf{Parameter} & \textbf{Description} & \textbf{Example} \\\hline
        bibliography\_entry\_key & \makecell{The key given to the entry when it was placed in \\\filename{references.bib}.} & biblatex\_manual \\
        \hline
    \end{newtable}
\end{command}

\begin{command}{Indirect Citation}
    An \underline{indirect citation} is used when referring to a bibliography entry as support for a statement, for example when saying \say{according to the paper (\textit{bibliography\_entry})...}. 
    Given the Harvard-style guidelines, this follows the general format:
    \begin{quote}
        \commandformat{(author, author, ..., and author, year\_published)}
    \end{quote}
    
    \noindent To call this function, we use the syntax given by \Cref{snippet:Indirect citation format.}.

    \begin{newsnippet}{Indirect citation format.}
        \begin{minted}{latex}
\indirectcitation{bibliography_entry_key}
        \end{minted}
    \end{newsnippet}

    \noindent The parameters for this function are shown below in \Cref{table:Indirect citation parameters.}.

    \begin{newtable}{Indirect citation parameters.}{|c|c|c|}
        \hline
        \textbf{Parameter} & \textbf{Description} & \textbf{Example} \\\hline
        bibliography\_entry\_key & \makecell{The key given to the entry when it was placed in \\\filename{references.bib}.} & biblatex\_manual \\
        \hline
    \end{newtable}
\end{command}

\newsection{Mathematics and Computer Science}

\comment{ADD THIS SECTION.}

\newsubsection{Additional Mathematics Commands}

\comment{ADD THIS SECTION.}

\newsubsection{Computer Science Environments}

\comment{ADD THIS SECTION.}